(1) 泡利矩阵满足 \(\{\sigma_y,\sigma_x\}=\{\sigma_y,\sigma_z\}=0\)，故 \(\{H_0,\sigma_y\}=0\)。记 \(y=x/\xi\)，由 \(\Delta\xi=\hbar v\)，零能方程等价于
\[
\partial_y\psi=-\tanh y\,\sigma_y\psi .
\]
令 \(\chi_\pm=(1,\pm i)^{\mathsf T}/\sqrt2\)，则 \(\sigma_y\chi_\pm=\pm\chi_\pm\)，各通道解为 \(\psi_\pm(y)=C_\pm(\cosh y)^{\mp1}\chi_\pm\)。只有正手征通道在两端均可归一化。由 \(\int_{-\infty}^{\infty}\operatorname{sech}^2y\,dy=2\)，
\[
\boxed{\psi_0(x)=\frac{\operatorname{sech}(x/\xi)}{\sqrt{2\xi}}\chi_+},\qquad
\boxed{H_0\psi_0=0,\quad\sigma_y\psi_0=\psi_0}.
\]
若质量剖面整体变号，零能一阶方程中的指数符号反转，唯一正规化通道改为 \(\chi_-\)。因此畴壁质量符号变化决定局域零模的手征，而非单凭势阱形状猜测。

(2) 用 \(\sigma_x\sigma_z=-i\sigma_y\) 严格展开平方，注意质量项与导数不对易：
\[
H_0^2=-\hbar^2v^2\partial_x^2+\Delta^2\tanh^2y
-\frac{\hbar v\Delta}{\xi}\operatorname{sech}^2y\,\sigma_y .
\]
定义 \(A=\partial_y+\tanh y\)，\(A^\dagger=-\partial_y+\tanh y\)（全直线上的形式伴随）。于是
\[
\frac{H_0^2}{\Delta^2}\Big|_{\sigma_y=+1}
=-\partial_y^2+1-2\operatorname{sech}^2y=A^\dagger A,\qquad
\frac{H_0^2}{\Delta^2}\Big|_{\sigma_y=-1}
=-\partial_y^2+1=AA^\dagger .
\]
负手征伙伴为自由算符加常数，其谱纯连续且从 \(1\) 起。若正手征通道存在非零本征值 \(0<\varepsilon^2<1\) 的平方可积态 \(f\)，则 \(Af\ne0\) 是 \(AA^\dagger\) 的同一离散本征值态，与自由伙伴无束缚态矛盾。零本征值只能满足 \(Af=0\)，唯一正规化解为 \(\operatorname{sech}y\)。故只有一条离散能级 \(E=0\)，连续谱为
\[
\boxed{E\in(-\infty,-\Delta]\cup[\Delta,\infty)}.
\]
上述结论依赖 \(\Delta\xi=\hbar v\)；不能未经证明推广到任意畴壁宽度。

(3) 对 \(k>0\)，自由伙伴 \(AA^\dagger\) 的平面波为 \(e^{iky}\)，用 \(A^\dagger\) 映射到正手征通道，得到
\[
f_k(y)=A^\dagger e^{iky}=(\tanh y-ik)e^{iky},\qquad
\bigl[-\partial_y^2+1-2\operatorname{sech}^2y\bigr]f_k=(1+k^2)f_k.
\]
它在 \(y\to-\infty\) 时为 \((-1-ik)e^{iky}\)，在 \(y\to+\infty\) 时为 \((1-ik)e^{iky}\)，完全没有反向波。因此左入射归一化透射振幅为
\[
\boxed{t(k)=\frac{1-ik}{-1-ik}=\frac{ik-1}{ik+1}},\qquad
\boxed{R=0,\quad T=|t(k)|^2=1}.
\]
对于选定的非零 \(E=\pm\Delta\sqrt{1+k^2}\)，完整自旋量可取
\[
\Psi_E(y)=C\left[f_k(y)\chi_+
+\frac{H_0}{E}\bigl(f_k(y)\chi_+\bigr)\right].
\]
因 \(H_0^2(f_k\chi_+)=E^2(f_k\chi_+)\)，直接有 \(H_0\Psi_E=E\Psi_E\)。两侧渐近质量的绝对值相同，入射与透射群速度大小相同；守恒概率流 \(j=v\Psi^\dagger\sigma_x\Psi\) 给出 \(R+T=1\)。阈值 \(k\to0^+\) 时 \(t\to-1\)，高能 \(k\to\infty\) 时 \(t\to1\)，两极限透射概率均为一；阈值相位不意味着反射。

(4) 由反对易关系，
\[
H_\delta^2=(H_0+\delta\sigma_y)^2=H_0^2+\delta^2\mathbb I.
\]
由于 \(\sigma_y\psi_0=\psi_0\)，原有正规化界面态不变且精确满足
\[
\boxed{H_\delta\psi_0=\delta\psi_0,\quad E_{\rm edge}=\delta}.
\]
连续谱能量满足 \(E^2\ge\Delta^2+\delta^2\)，其正负边缘为
\[
\boxed{E_{\rm cont,\pm}=\pm\sqrt{\Delta^2+\delta^2}}.
\]
因为 \(|\delta|<\sqrt{\Delta^2+\delta^2}\) 对任意 \(\Delta>0\) 严格成立，界面态始终处在体能隙中且保持相同局域长度。可是 \(\delta\sigma_y\) 不再与 \(\sigma_y\) 反对易，零能谱对称性被破坏，故界面态虽局域却不再钉扎于零能。这说明零能保护与束缚态存在性是不同命题。